%% chapters/behavior_modeling/ch2_correlation.tex
%% 第2章　相関・異質性
\section{MNLからの拡張: 相関・異質性}\label{sec:abms}
本節では，前節で導入したMNLモデルからの拡張として，選択肢間の相関や意思決定主体間の異質性を考慮したモデルを解説する．
具体的には，選択肢間の相関を部分的に許すNested Logitモデル，個人の異質性を考慮するMixed Logitモデル，そして経路選択特有の拡張であるPath-Size Logitモデルなどを紹介する．

\subsection{IIAの限界とNested Logitモデルの導入}
\subsubsection{IIA特性と赤バス・青バス問題}
IIA（Independence of Irrelevant Alternatives）特性は直訳すると「無関係な選択肢との独立性」となるが，具体的には任意の二つの選択肢の間の選択確率の比は、その二つの選択肢の効用のみによって決まるという性質のことである．
例えば選択肢 $a$ と $b$ の間の選択確率の比は，以下の様に選択肢 $a$ と $b$ の効用のみによって決まる．
\begin{equation}
  \frac{P(a)}{P(b)} = \frac{\frac{\exp(V_a)}{Z}}{\frac{\exp(V_b)}{Z}} = \exp(V_a - V_b)
\end{equation}
ただし $Z = \sum_{a'} \exp(V_{a'})$ である．
つまり選択肢 $a$, $b$ いづれとも異なる第三の選択肢 $c$ の効用 $V_c$ はこの比に影響を与えない．
このIIA特性は，各選択肢効用の誤差項が独立であるとするi.i.d.の仮定に起因するものであり，MNLモデルの重要な特徴である．

しかしながら選択肢間の相関が存在する場合にはこのIIA特性が崩れる．
その代表的な例として\textbf{赤バス・青バス問題}が知られている．

\begin{itembox}[l]{赤バス青バス問題}
    「車」「赤バス」の二つの選択肢から選ぶ二項ロジットモデルにおいて，効用の確定部分が全く同じであれば両者の選択確率は $1/2, 1/2$ となる．
    ここに「青バス」という選択肢を加えたとき，これの効用の確定部分も他と同じであったとすれば三つの選択肢の選択確率は $1/3, 1/3, 1/3$ となる．
    色違いのバスが登場しただけでもともと車を利用していた人の三分の一もがバスに流れるのは不自然である．
\end{itembox}
赤バス青バス問題は，車と赤バスの選択確率の比が青バスの有無によらず $1:1$ であるというIIA特性によって生じる問題である．
より根本的には，赤バスと青バスの効用の誤差項が独立であるというi.i.d.の仮定が現実には成立しないことに起因する．
すなわち，（「色違いのバスが登場しただけでもともと車を利用していた人の三分の一もがバスに流れるのは不自然である」という主張が暗に仮定するように）赤バス・青バスはともに車と比べて公共交通という性質を共有している点で類似した選択肢であるため，両者の効用の誤差項は相関していると考えるのが自然である．
こうした選択肢間の相関を考慮するためには，\textbf{Nested Logit モデル}などのより応用的なモデルを用いる必要がある．


\subsubsection{Nested Logitモデルの定式化}
Ben-Akiva(1973

\subsection{経路相関とPath Size Logitモデル}


\subsection{個人間の異質性とMixed Logit モデルの導入}

% 第3節 MNLからの拡張：相関・異質性

% ここで学生が最初につまずくので、拡張理由を明確に書くべきです。

% 3.1 IIAとその限界
% 類似選択肢があると不自然
% 経路の重複問題
% 3.2 Nested Logit
% 誤差相関を部分的に許す
% 木構造による整理
% 3.3 Mixed Logit
% 係数異質性
% 柔軟性
% シミュレーション推定
% 3.4 Path-Size Logit など経路選択特有の拡張
% 経路重複への対応
% 経路選択の相関は交通分野で特有に重要

% 3.5
% 観測異質性と非観測異質性
% 交互作用項で表せること／表せないこと

% ここで大事なのは、
% Nested Logit は「選択肢相関」への対応、Mixed Logit は「個人異質性 + 柔軟な相関」への対応
% と役割を分けて説明することです
